\section*{Techniques from Hilbert Space Theory}

\subsection*{12.1 L2-Norm and Inner Product Space}

This chapter presents a geometric perspective on random variables by viewing them

as vectors in a vector space. Although this vector space has infinite dimension

in general, this viewpoint offers insights into the behavior of random variables.

Assuming that all random variables in the vector space have finite second moments,

we can define an inner product that resembles the dot product in a finite-dimensional

Euclidean space. This inner product allows us to define notions such as orthogonal-

ity and the analog of the triangle inequality in this infinite-dimensional vector space.

To study random variables form this geometric perspective, we borrow tech-

niques from Hilbert space theory. We can define the projection operator, which

maps a random variable to a subspace, and prove the Orthogonality Principle in

the context of probability theory. These tools are then applied to the derivation of

statistical estimators that minimize mean-squared error.

\begin{defbox}{12.1}
\end{defbox}

Given a measure space (\Omega, \mathcal{F},\mu ), we denote the set of all \mathcal{F}-measurable

functions X with finite second moment by L2(\mu)More generally, for any

positive integer p, we define Lp (\mu) as the set of all \mathcal{F}-measurable functions

with finite

\int

\vertX\vertp d\mu.

In this chapter, we will focus on the study of L2(P) for a probability measure P .

We note that any square-integrable random variable necessarily has a finite mean,

implying that L2(P) is a subset of L1(P) Using Cauchy-Schwarz inequality

(Theorem 11.1), we can define an inner product for two square-integrable random

variables.

\begin{defbox}{12.2}
\end{defbox}

Let (\Omega, \mathcal{F},P) be a probability space. We define the inner product of two real-

valued random variables in L2(P) by

\langleX, Y\rangle\coloneqqE[XY ].

We define the L2-norm of X by

X2 \coloneqq

\sqrt

\langleX, X\rangle=

\sqrt

E[X2].

Two random variables X and Y are said to be orthogonal if \langleX, Y\rangle= 0. We will

sometime write Xwithout the subscript 2 to simplify notation.

This notion of orthogonality is analogous to the concept of perpendicularity in

Euclidean geometry. When random variables X and Y have zero mean, then they

are orthogonal iff they are uncorrelated.

Complex Inner Product For complex random variables, the inner product is

defined as E [X*Y], where X* denotes the complex conjugate of X.

\textbf{Example 12.1.1 (A Finite Inner Product Space)} \\

Consider a finite probability space with sample space \Omega ={ a, b, c}The \sigma-algebra is the discrete

\sigma-algebra P(\Omega) , and the probability measure is given by P({a}) = pa , P({b}) = pb, P({c}) =

pc,w h e r e pa , pb,a n d pc are nonnegative real numbers that sum to 1.

We can specify a random variable X by giving its value at the three points in \Omega , which can be

represented as a vector (X(a), X(b), X(c))Conversely, any three-dimensional vector (x1,x 2,x 3)

is associated to a random variable X by defining X(a) \coloneqqx1, X(b) \coloneqqx2,a n d X(c) \coloneqqx3. Thus,

the random variables in this probability space have a one-to-one correspondence with the vectors

in a three-dimensional vector space over R.

The inner product between two random variables X and Y , which are represented by vectors

(x1,x 2,x 3) and (y1,y 2,y 3), respectively, is computed as

\langleX, Y\rangle= pax1y1 + pbx2y2 + pcx3y3.

The second moment of a random variable X is thus equal to pax2

1 + pbx2

2 + pbx2

3 .

\textbf{Example 12.1.2 (A Complex Hilbert Space of Infinite Dimension)} \\

Let \Omega ={ 1, 2, 3,... } be the sample space, and the power set P(\Omega) be the \sigma-algebra. Consider a

probability measure P defined on \Omega by

P({k}) \coloneqq 1

2k

for k \geq 1. A random variable X can be identified with an infinite complex sequence

X (x1,x 2,x 3,... ) ,

where xk are complex numbers for k = 1, 2, 3,... .

The random variable X belongs to L 2(P) if and only if the series \sum\infty

k=1 \vertxk\vert2/2k is convergent.

If we are given another random variable Y , whose values at 1 , 2, 3, 4,.are y 1,y 2,y 3,y 4,... ,t h e

inner product between X and Y is given by

\langleX, Y\rangle=

\infty\sum

k=1

x*

k yk/2k.

Here, x *

k denotes the complex conjugate of x k .

The followings are other basic properties of inner product.

\begin{thmbox}{12.1}
\end{thmbox}

Let\langle\cdot,\cdot\ranglebe an inner product over a field K , where K is either R or C :

Fo r X, Y1,Y 2 \in L2(P) and \alpha, \beta \in K,

\langleX, \alphaY1 +\betaY 2\rangle= \alpha\langleX, Y1\rangle+\beta\langleX, Y2\rangle.

For real inner product, we have\langleX, Y\rangle=\langle Y, X\ranglefor X, Y \inL 2(P) .

For complex inner product, we have\langleX, Y\rangle=\langle Y, X\rangle* for X, Y \in L2(P) .

Fo r X \in L2(P) and \alpha \inK ,\alphaX2 =\vert \alpha\vert\cdot X2.

The triangle inequality for L 2 norm is called the Minkowski inequality.

\begin{thmbox}{12.2 (Minkowski Inequality)}
\end{thmbox}

For any two random variables X and Y in L 2(P) ,

X+ Y 2 \leq X2 + Y2.

\textit{Proof} We give a proof for complex-valued random variables. Using the Cauchy-

Schwarz inequality (Theorem 11.1),

X+ Y 2

2 \leqE [\vertX\vert2]+ 2R e(E[X*Y])+ E [\vertY\vert 2]

\leqE [\vertX\vert2]+ 2

\sqrt

E[\vertX\vert2]E[\vertY\vert 2]+ E[\vertY\vert 2]

=( X2 + Y2)2.

Taking the square roots of both sides yields the Minkowski inequality. \hfill $\square$

Similar to the case of L1 norm, a random vector with zero L2 norm is equivalent

to the zero function but need not be equal to the zero function. We continue with the

convention that two random variables that are equal almost everywhere are regarded

as the same random variable. With this caveat, we have shown that L2(P) is a

normed vector space.

Using the Minkowski inequality, we can see that the L2 norm defines a metric

on the space L2(P) , by regarding X - Y2 as the distance between two random

variables. Moreover, we can prove that the space of random variables in L2(P) is

complete with respect to the norm induced by inner product. It is a consequence of

the Riesz-Fischer theorem.

\begin{defbox}{12.3}
\end{defbox}

A sequence of vectors (un)\infty

n=1 in a vector space V with a norm function \cdot is

called a Cauchy sequence if given any \epsilon> 0, there exists an integer N such that

um - un\leq \epsilon for all m, n \geq N. We say that a sequence of vectors (un)\infty

n=1 is

convergent if there exists a vector u0 \in V such that for any given \epsilon> 0, there is

an integer N such that un - u0\leq \epsilonfor all n \geq N .

A normed vector space V is said to be complete if every Cauchy sequence in V is

convergent. A Hilbert space is a (real or complex) vector space that is complete

with respect to the norm induced by the inner product.

\begin{thmbox}{12.3 (Riesz-Fischer)}
\end{thmbox}

Let (\Omega, \mathcal{F},P) be a probability space. The space L2(P) is a complete vector

space with respect to the L2 norm.

\textit{Proof} Suppose we have a Cauchy sequence f1, f2, f3,.in L2(P) To prove

that the sequence converges in mean square, we need to find a candidate function

and show that it is indeed the limit. Our first step is to show that there exists a

subsequence that converges almost surely.

By the definition of Cauchy sequence, we can find a sequence of integers n1 <

n2 <n 3 < \cdot\cdot\cdot such that fnk+1 - fnk 2 \leq 1/2k for k = 1, 2, 3,... We claim that

the limit limk\to\infty fnk () exists almost everywhere.

We write fnk () as

fnk () = fn1 () +

k- 1\sum

i=1

(fni+1() - fni ())

and apply the result established in Problem 7.7, which implies that the limit of

fnk () exists almost everywhere if

\int

\vertfn1\vertdP +

\infty\sum

i=1

\int

\vertfni+1 -f ni \vertdP (12.1)

is finite.

To show that it is finite, we use Jensen inequality to obtain

\int

\vertfni+1 -f ni \vertdP \leq

( \int

\vertfni+1 -f ni \vert2 dP

) 1/2

= fni+1 -f ni 2

for all i Our choice of the indices n i 's ensures thatfni+1 -fni 2 \leq1 /2i Therefore,

the sum in (12.1) is finite. By applying the result from Problem 7.7, there exists an

event E with probability 1 such that f nk () converges for all \in E. We can define

a candidate function as follows:

f( )\coloneqq

{

limk\to\infty fnk ()for \in E,

0 otherwise .

If is not in E ,t h e v a l u eo f f( ) is arbitrary. We set the value of f( ) to be 0 for

\in Ec. This convention ensures that f( ) is measurable.

Let\epsilon> 0 be an arbitrarily small real number. By the choice of the sequence n j ,

there is a sufficiently large integer m such thatfnj - fnk 2

2 \leq\epsilon for all j, k \geq m.

We fix k \geq m and let j \to\infty Then, by applying Fatou's lemma (Theorem 7.3), we

have

\int

E

lim

j\to\infty

\vertfnj -f nk \vert2 dP\leq lim

j\to\infty

\int

E

\vertfnj -f nk \vert2 dP\leq \epsilon.

Since f nj () \to f( ) for \in E, we get f - fnk 2

2 =

\int

\vertf - fnk \vert2 dP \leq\epsilonfor

all k \geq m. Therefore, the subsequence (f nk )k\geq 1 converges to f in mean square.

To show that the original sequence (f n)\infty

n=1 also converges to f in mean square,

we write

fn -f 2 <fn -f n\ell2 + fn\ell -f 2,

and take limit as n and\ell approach infinity. Since (f n)n\geq1 is a Cauchy sequence,

the termfn -f n\ell2 tends to zero. Furthermore, from the previous paragraph, we

know thatfn\ell -f2 tends to zero as well. This proves thatfn -f2 can be made

arbitrarily small for all sufficiently large integer n , and thus (f n)n\geq1 converges to f

in mean square. Because the L 2 norm is a continuous function (Exercise 12.2), the

L2 norm of f is finite. This proves that f \in L2(P) \hfill $\square$
